%! Comparing Distributions — Unit 2, Lesson 2.4 — Homework (Answer Key)
% This is the lesson's SCORED individual practice and the LAST component in the packet.
% Unit 2 on (2026-09-29): ASSIGNED at Close & Assign and done OUTSIDE class — not started in
% class. Packet pages are the default; a Desmos activity may replace them on a given day
% (decided at Close & Assign).
% THIRD CONTEXT: the notes teach on pizza delivery times (Tony's, Rosa's, Slice House), the
% Guided Practice runs on paper airplanes, and these items are on phone batteries.
% TWO PAGES MAX: two boxes — items 1-2 on page 1, items 3-6 and the spiralbox on page 2 — so no
% item breaks across a page.
% Items 4 and 5 are the deliberate CONTRAST PAIR: the same question with opposite needs, two
% brands with the SAME median — both answers are decided by spread, and they point to
% opposite brands. Item 6 is the crux head-on (a centers-only claim). Item 3 is the spiral
% (Lessons 2.2-2.3: the mean is not resistant, the median is).
% THE DATA (quartiles: median of each half, overall median left out when n is odd)
%   Volt (n=11): 6 10 11 11 12 12 12 13 13 14 15
%     min 6, Q1 11, med 12, Q3 13, max 15; IQR 2; fences 8 and 16 -> 6 is an outlier;
%     mean 129/11 = 11.7; without the 6: mean 123/10 = 12.3, median (12+12)/2 = 12
%   Nova (n=11): 4 7 8 10 11 12 14 15 16 18 19
%     min 4, Q1 8, med 12, Q3 16, max 19; IQR 8; fences -4 and 28 -> no outliers
%   at least 10 h: Volt 10/11, Nova 8/11;  at least 16 h: Volt 0/11, Nova 3/11
% PAGE LOCKSTEP: the key differs only in -key for -boxes, the header, \blank -> \ans and
% \writelines{n} -> n \ansline; every work block is byte-identical.
\documentclass[12pt]{article}
\usepackage{saar-article}
\usepackage{saar-key}   % pulls in -boxes; do NOT also load -boxes
\newcommand{\TallMath}[1]{$\displaystyle #1\rule[-1.4em]{0pt}{3.2em}$}
% Word bank strip for every fill-in cluster. Defined per document; the key inherits it and
% prints the same strip, so heights match.
\newcommand{\wordbank}[1]{\par\vspace{2pt}\noindent\fcolorbox{goldacc}{goldbg}{\parbox{\dimexpr\linewidth-2\fboxsep-2\fboxrule\relax}{\small\textbf{Word bank:} #1}}\par\vspace{4pt}}
% Handwriting room inside every \begin{work} block. Moves the blank and the
% key together, so the two can never drift apart.
\setlength{\workrowsep}{10pt}

\begin{document}

\pageheader{Unit 2, Lesson 2.4}{Homework --- Answer Key}
% NAMESTRIP: no \namedateperiod here — mirrors the blank. NO teachernote here —
% teacher prose lives in the lesson plan.

\vspace{0.12in}

% Authored identically in homework/ and homework_key/.
\begin{remindbox}
\small
\textbf{This is your graded homework.} It is scored out of 12 (2 points per item) and due at
the start of the first class after two study halls. It is assigned today and done on your own,
outside class --- flip back to the \emph{Keep in Mind} box on the cover and to rows 3 and 4 of
your notes.
\end{remindbox}

\vspace{0.08in}

\boxguard[20]
\begin{scenariobox}[Volt vs.\ Nova Phone Batteries]{cerulean}
A tech club fully charged $11$ phones of each brand and played the same video on every phone
until its battery died. Battery life, in hours, in order:

\vspace{3pt}
\textbf{Volt:} \quad $6,\ 10,\ 11,\ 11,\ 12,\ 12,\ 12,\ 13,\ 13,\ 14,\ 15$
\par\vspace{2pt}
\textbf{Nova:} \quad $4,\ 7,\ 8,\ 10,\ 11,\ 12,\ 14,\ 15,\ 16,\ 18,\ 19$
\end{scenariobox}

\vspace{0.08in}

\boxguard
\begin{notesbox}{Comparing Two Brands}
Every answer names the brand and the hours.
\wordbank{4 \;\textbullet\; 6 \;\textbullet\; 8 \;\textbullet\; 11 \;\textbullet\; 12 \;\textbullet\; 13 \;\textbullet\; 15 \;\textbullet\; 16 \;\textbullet\; 19 \;\textbullet\; 2 \;\textbullet\; outlier}
\begin{enumerate}[leftmargin=1.6em, itemsep=9pt, topsep=3pt]

  \item Complete each brand's five-number summary and IQR.

        \begin{center}
        \renewcommand{\arraystretch}{1.8}
        \begin{tabularx}{\linewidth}{@{} l Y Y Y Y Y Y @{}}
        \toprule
        \textbf{Brand} & \textbf{Min} & $\boldsymbol{Q_1}$ & \textbf{Median} & $\boldsymbol{Q_3}$ & \textbf{Max} & \textbf{IQR} \\
        \midrule
        Volt & \ans{6} & \ans{11} & \ans{12} & \ans{13} & \ans{15} & \ans{2} \\
        Nova & \ans{4} & \ans{8} & \ans{12} & \ans{16} & \ans{19} & \ans{8} \\
        \bottomrule
        \end{tabularx}
        \end{center}

        In context, what does Nova's IQR tell you?
        \par\ansline{Middle half of Nova phones: 8 to 16 hours, an 8-hour span.}

  \item Is Volt's $6$-hour phone an outlier? Show the lower fence,
        $Q_1 - 1.5 \cdot \text{IQR}$.

        \begin{work}
          \text{Lower fence} &= 11 - 1.5(2) \\
                             &= 8 \text{ hours}
        \end{work}

        \vspace{4pt}
        Outlier? \ans{Yes --- 6 hours is below the lower fence of 8 hours.}

\end{enumerate}
\end{notesbox}

\boxguard[30]
\begin{notesbox}{Comparing Two Brands, continued}
\begin{enumerate}[leftmargin=1.6em, itemsep=5pt, topsep=2pt, start=3]

  \item \emph{Back to Lessons 2.2--2.3.} Find Volt's mean with the $6$-hour phone and without
        it, and Volt's median both ways. Then: which pair of measures --- mean and SD, or
        median and IQR --- should you use to compare these brands, and why?

        \begin{work}
          \bar{x}_{\text{with the 6}} &= 129 / 11 \approx 11.7 \text{ hours} \\
          \bar{x}_{\text{without it}} &= 123 / 10 = 12.3 \text{ hours} \\
          \text{median} &= 12 \text{ hours both ways}
        \end{work}

        \par\ansline{Median and IQR: Volt's 6-hour outlier pulls the mean}
\ansline{(11.7 vs.\ 12.3 hours) but not the median, which stays 12.}

  \item Maya needs a phone that lasts a full school day: \textbf{at least $10$ hours}. How
        many of each brand's $11$ phones did that? Which brand should she buy?
        \par\ansline{Volt: 10 of 11. Nova: 8 of 11 (the 4, 7, and 8 fall short).}
\ansline{Maya should buy Volt --- same median, more consistent.}

  \item Leo wants his best chance at a phone that lasts \textbf{at least $16$ hours} on a
        long trip. How many of each brand's $11$ phones did that? Which brand should he buy?
        \par\ansline{Volt: 0 of 11. Nova: 3 of 11 (16, 18, and 19 hours).}
\ansline{Leo should buy Nova --- its greater spread reaches higher.}

  \item Jordan says, \emph{``Both brands have a median of $12$ hours, so it doesn't matter
        which one you buy.''} Say what Jordan left out, then write the comparison with
        comparative words, in context.
        \par\ansline{Jordan compared centers only. Both medians are 12 hours, but}
\ansline{Nova's IQR (8 h) is 4 times Volt's (2 h), so Volt's battery}
\ansline{life is far more consistent. Volt has one low outlier (6 h).}

\end{enumerate}
\end{notesbox}

\boxguard[5]   % the short spiralbox needs about five baselines; a larger guard risks a third page
\begin{spiralbox}
\textbf{Next class: Unit 2 Review.} This was the last lesson of Unit~2. Next class we review
for the Unit~2 test: displays and SOCS, mean, median, SD and IQR, boxplots and fences, and
today's comparisons. Item~6 is exactly the kind of question the test asks.
\end{spiralbox}

\end{document}
